% Specific Aims — NIH R01 (standalone, 1 page, no bibliography)
% Compiled with: lualatex
\documentclass[11pt,letterpaper]{article}

% -- Fonts (NIH: >= 11pt, Arial) --
\usepackage{fontspec}
\setmainfont{Arial}

% -- Page geometry (NIH: >= 0.5in margins) --
\usepackage[margin=0.5in]{geometry}
\usepackage{setspace}
\singlespacing

% -- Formatting --
\usepackage[normalem]{ulem}  % \uline for underline without changing \emph
\usepackage{enumitem}

% -- Placeholder text (remove with ./strip-lipsum.sh once you start writing) --
\usepackage{lipsum}


% -- No paragraph indent, compact skip --
\setlength{\parindent}{0pt}
\setlength{\parskip}{0pt}

% -- No page numbers --
\pagestyle{empty}

\begin{document}
\enlargethispage{0.4in}
\vspace*{-0.15in}
\textbf{SPECIFIC AIMS}
\vspace{-2pt}

% TODO: Opening paragraph — set the scientific context. Two to four sentences
%   establishing the field, the unmet need, and why this problem matters now.
\lipsum[1]

\vspace{2pt}
% TODO: Problem paragraph — frame the critical gap. Make the limitation of
%   current approaches concrete; end with a sentence that motivates the proposal.
\lipsum[2]

% TODO: Long-term goal / objective paragraph. One or two sentences on the
%   long-term goal of the research program and what this proposal
%   accomplishes. A short transition sentence introduces the aims.
\lipsum[3]

\textbf{Aim 1: <Aim 1 title>.}
% TODO: 2--4 sentences describing Aim 1.
\textit{Aim 1A:\ <Subaim 1A title>.}
% TODO: 1--2 sentences on Subaim 1A.
\textit{Aim 1B:\ <Subaim 1B title>.}
% TODO: 1--2 sentences on Subaim 1B.

\textbf{Aim 2: <Aim 2 title>.}
% TODO: 2--4 sentences describing Aim 2.
\textit{Aim 2A:\ <Subaim 2A title>.}
% TODO
\textit{Aim 2B:\ <Subaim 2B title>.}
% TODO

% -- Optional Aim 3 (uncomment for 3-aim grants) --
% \textbf{Aim 3: <Aim 3 title>.}
% % TODO: 2--4 sentences describing Aim 3.
% \textit{Aim 3A:\ <Subaim 3A title>.}
% % TODO
% \textit{Aim 3B:\ <Subaim 3B title>.}
% % TODO

% TODO: Closing impact statement (one sentence in bold) summarizing what the
%   aims together will deliver.
\textbf{Together, these aims will % TODO
.}

\end{document}
